


\textbf{Update of the delta verification run (2026-11-16).} The byte layout that this check found missing is now specified in \texttt{docs/05\_architecture/FRAME\_LAYOUT.md}, and its section 8 resolves each mismatch below (five mismatches and one gap) by a document change or a stated abstraction, with ten further rows. The tables below are the record of the check as it was made.

Method. The message terms of every model file in \texttt{formal/} were read next to the byte layouts in \texttt{src/crashguard/packet.py}, \texttt{src/crashguard/sct\_protocol.py} and \texttt{firmware\_rd/cg\_rd.h} and \texttt{cg\_rd.c}. The paper has no frame-layout figure; frame sizes appear in the text and in the architecture figure (\texttt{paper/figures/pp\_arch.png}): 97-byte alert, 26-byte acknowledgement, 14-byte flash and cancel, 8-byte tag. Those sizes agree with the code (table 1).

\subsection*{1. Sizes: code against paper}

{\scriptsize\raggedright
\begin{longtable}{p{0.190\linewidth}p{0.190\linewidth}p{0.190\linewidth}p{0.190\linewidth}p{0.190\linewidth}}
\toprule
Frame & Python reference & C firmware path & Paper & Match \\
\midrule
\endhead
Full alert & 97 bytes: 89 body + 8 tag (\texttt{ALERT\_LEN}, \texttt{ALERT\_BODY\_LEN}, \texttt{TAG\_LEN}) & 97 / 89 / 8 (\texttt{CG\_ALERT\_LEN}, \texttt{CG\_ALERT\_BODY}, \texttt{CG\_TAG\_LEN}); \texttt{cg\_encode\_alert} writes the same field order & 97 bytes & yes \\
ACK & 26 bytes: 18 body + 8 tag & 26 / 18 & 26 bytes & yes \\
Provisional flash and cancel & 14 bytes: 6 body + 8 tag (\texttt{FLASH\_LEN}, \texttt{FLASH\_BODY\_LEN}) & no flash encoder in \texttt{cg\_rd.c} (only the D1 retry logic, the alert and the freshness ring) & 14 bytes & yes (Python only; the flash is not in the C path) \\
Tag & HMAC-SHA256 truncated to 8 bytes over all preceding body bytes & the same (\texttt{cg\_hmac\_sha256}, 8 bytes, constant-time compare) & 8 bytes; 2\textasciicircum{}-64 per forgery attempt (ANALYTICAL) & yes \\
\bottomrule
\end{longtable}}

\subsection*{2. Fields: models against the layout}

{\scriptsize\raggedright
\begin{longtable}{p{0.158\linewidth}p{0.158\linewidth}p{0.158\linewidth}p{0.158\linewidth}p{0.158\linewidth}p{0.158\linewidth}}
\toprule
Item & Layout (reference and C path) & Tamarin \texttt{crash\_alert\_auth} (M1 and M0) & Tamarin \texttt{crash\_alert\_fresh} and \texttt{\_bootctr} & ProVerif \texttt{carid\_*} & Result \\
\midrule
\endhead
Alert fields & version, type, node\_id (4), event\_id (4), counter (4), timestamp\_ms (4), 12 payload fields, 47 reserved, tag (8) & \texttt{<'ALERT', v, c, cls, mac(key, <'A', v, c, cls>)>}: vehicle, one fresh field, class & same with field \texttt{<e, n>} (epoch, sequence) or \texttt{<b, e, n>} & \texttt{(v, id, mac(kdf(km, v), (v, id)))} & abstraction: every payload field is one value \texttt{cls}; event\_id, counter and timestamp are one name. Acceptable for the lemmas (the tag covers all bytes). No mismatch of the properties proved \\
ACK fields & version, type = 2, node\_id, acked event\_id, acked counter, base counter, tag & \texttt{<'ACK', v, f, mac(kdf(km, v), <'ACK', v, f>)>} for M1; shared key \texttt{kdf(km, 'cell')} for M0 & same ACK & not modelled (car side only) & the base counter is not modelled (it does not enter any lemma); \texttt{f} stands for the pair (event\_id, counter). M0 and M1 differ as the paper says (key and binding) \\
Per-vehicle key & K\_i = HMAC(master, vehicle id), one KDF per new id, cached (\texttt{ar\_r4.kdf}) & \texttt{kdf(km, v)} & \texttt{kdf(km, v)} & \texttt{kdf(km, v)} & match (a symbolic function stands for the HMAC) \\
Domain-separation letters & message type inside the tagged body (alert 1, ACK 2, flash 4, cancel 5) and different lengths & letters \texttt{'A'} (alert) and \texttt{'ACK'} & \texttt{'A'} & none (one message kind) & match in effect for the alert and the ACK. The ProVerif car-side models have no letter; they model one kind of frame, so no lemma is affected. The round-4 model \texttt{crash\_alert\_prov} adds \texttt{'P'}; its control \texttt{\_nodomain} (one letter for both) is falsified as expected \\
Short node id & the flash carries only the low 16 bits of node\_id (\texttt{sct\_protocol}); the alert carries 32 bits & vehicle \texttt{v} is one name & same & same & \textbf{mismatch (assumption)}: uniqueness of the 16-bit id within the fleet is an assumption of the layout and is not part of any model \\
Freshness field & \texttt{PACKET\_FORMAT.md} and \texttt{sct\_protocol.py}: a persisted 32-bit counter (the flash carries its low 16 bits, rebuilt within +-32,768 of the last value the base has seen). \texttt{ar\_r1.py}, \texttt{ARCHITECTURE.md} (design R1) and the C ring \texttt{cg\_fresh\_accept(field, seq, ...)}: field = GPS seconds, sequence in second, no persisted counter & a fresh name \texttt{c} (no ordering) & id \texttt{<e, n>}: epoch = GPS second, \texttt{n} = 1 to 3 in that second & \texttt{mk(e, s1)} & \textbf{mismatch (documentation)}: the byte positions of the GPS-second field and of the sequence-in-second are not specified in \texttt{PACKET\_FORMAT.md}, which still describes the persisted counter. The freshness models and \texttt{ar\_r1} describe a different frame semantics from \texttt{PACKET\_FORMAT.md} \\
Window & C path: \texttt{cg\_fresh\_accept(..., W, R)}; the simulation uses W = 10 s (\texttt{ar.r1}) & no window & window of two epochs (the current and the previous GPS second) & no window & \textbf{mismatch (abstraction)}: the model's window (about 1 to 2 s) is not the simulated W = 10 s. The lemma (at most once per window) does not depend on the width, but the statement "within the window" is for the model's window \\
Seen-set & ring of the last R entries (R up to 512), cleared at reboot & none & unbounded set & none & \textbf{mismatch (abstraction, stated in the model file)}: a finite ring is outside the model. A replay of an entry that the ring has dropped is possible for a window longer than the ring covers (\texttt{ar.r1} measures the replay rate inside the window; H-R1b was not supported at the held-out burst load) \\
Boot counter & not in any frame layout (the alert has 47 reserved bytes; the flash has none) & none & fresh value per boot in the variant & fresh value per boot in \texttt{carid\_bootctr\_*} & \textbf{mismatch}: the paper says the variant costs one byte per frame. The alert has room. The flash would grow from 14 to 15 bytes. Airtime does not change: 13, 14 and 15 bytes take 46.34 ms at SF7 and 164.86 ms at SF9 (\texttt{lora.airtime\_ms}), so the airtime results hold. The layout change is not specified anywhere \\
14-byte flash and cancel & \texttt{sct\_protocol.py} & not in any model of the frozen design & not in any model & not modelled & \textbf{gap}: the flash and the cancel frames of the frozen design have no formal model. Only the 24-byte provisional frame of round 4 (\texttt{crash\_alert\_prov*.spthy}) was modelled, and that design was not supported (H-A1a, H-A1b). The flash and the full alert of one event carry the same counter in the layout; the models give each frame its own fresh id \\
Tag length & 8 bytes & perfect MAC (no truncation) & perfect MAC & perfect MAC & stated limit: symbolic model with perfect cryptography; the 2\textasciicircum{}-64 forgery bound is separate and ANALYTICAL \\
\bottomrule
\end{longtable}}

\subsection*{3. What the paper may say}

\begin{itemize}
\item The Tamarin results apply to the alert and ACK exchange of the full 97-byte alert and 26-byte ACK, abstracted to (vehicle, freshness id, class, tag).
\item No formal result is claimed for the 14-byte flash, the cancel frame, the 16-bit short node id or the finite ring.
\item The freshness lemmas are for the GPS-second id of the freshness design, not for the persisted counter of \texttt{PACKET\_FORMAT.md}; the byte layout of that design is not yet specified.
\item The boot-counter variant needs a layout change in the flash frame (14 to 15 bytes, airtime unchanged at SF7 and SF9).
\end{itemize}

These points are carried into the paper (limits paragraph of the formal section) and into \texttt{paper/PLACEHOLDERS.md} where the author must decide the byte layout.

